\documentclass[11pt]{amsart}
\usepackage[T1]{fontenc}
\usepackage[utf8]{inputenc}
\usepackage{amssymb}
\usepackage[margin=1.05in]{geometry}
\usepackage[hidelinks]{hyperref}
\newcommand{\stamptime}{\the\numexpr\time/60\relax:\ifnum\numexpr\time-60*(\time/60)\relax<10 0\fi\the\numexpr\time-60*(\time/60)\relax}
\AddToHook{shipout/foreground}{\put(0.5\paperwidth,-\paperheight+0.55cm){\makebox[0pt]{\tiny\ttfamily unit/doc/\jobname.tex \ --- \ compiled \today, \stamptime}}}
\newtheorem{theorem}{Theorem}
\newtheorem{proposition}[theorem]{Proposition}
\newtheorem{lemma}[theorem]{Lemma}
\newtheorem{corollary}[theorem]{Corollary}
\theoremstyle{definition}
\newtheorem{definition}[theorem]{Definition}
\theoremstyle{remark}
\newtheorem{remark}[theorem]{Remark}
\newcommand{\Om}{\Omega}
\newcommand{\Map}{\mathrm{Map}_*}
\newcommand{\Z}{\mathbb{Z}}
\newcommand{\code}[1]{\texttt{#1}}
\newcommand{\cT}{\mathsf{isTrunc}}
\newcommand{\cC}{\mathsf{isConn}}
\newcommand{\refl}{\mathsf{refl}}
\newcommand{\ap}{\mathsf{ap}}
\newcommand{\Hom}{\mathrm{Hom}}
\newcommand{\hc}[2]{[#1,#2]_*}
\newcommand{\Fib}{\mathrm{Fib}}
\newcommand{\tr}[1]{\|#1\|}
\newcommand{\cls}[1]{[#1]}
\title{The mapping theorem for targets $\Fib k$}
\date{28 October 2026}
\begin{document}
\maketitle

\noindent\textbf{Purpose.} Informal proof for \S2 of
\code{doc/pi3-s2-postnikov-handoff-7.txt} (\S3c of the plan): for
$P=\Fib\kappa$ and every Postnikov triple $(A,B,k)$, the canonical
projection $q_*:\hc P{\Fib k}\to\hc P{K(A,2)}$ is a bijection; and,
for a class $s\in\hc{S^2}P$ constructed here, precomposition
$s^*:\hc P{\Fib k}\to\hc{S^2}{\Fib k}$ is a bijection. Written before the
formal proof; numbered statements are meant to be declarations.
Foundations as in the development (univalence, images; function
extensionality derived); no new axiom; the sphere axiom
(\code{SphereAxiom}) is used in \S\S4--5 only.

\smallskip\noindent\textbf{Hypotheses.} $\kappa:K(\Z,2)\to_*K(\Z,4)$ is
any pointed map with (H4) (\code{KZ2Cohomology.H4}); we use the
\emph{surjectivity} of $\Phi_B$ (\S3) and, through
\code{FibreH3.FibAt[$\kappa$,h4k,B].merely-zero}, its injectivity ((H3)
enters there through \code{H3Proof.h3}, not as a hypothesis).
Instantiated with \code{CupSquare.kappa}, \code{H4Proof.h4}.
\emph{Refereed} independently (a fresh reader given the note and the
library, not the conclusion; and the user): correct, with
formalisation-level fixes now incorporated ($e_1^{-1}$ is the map of the
pointed equivalence \code{pt-inv e1}; the cancellation in
Proposition~\ref{p:right} is on the other side, \code{cancel-r}; (U) for
$P$ is \code{merely-zero}; the connectivity of $T_x$ inlined; the wording
on lift spaces).

\smallskip\noindent\textbf{The idea.} One lifting lemma (\S2), for an
arbitrary pointed source $Q$: if (E) every composite $F;k$ with
$F:Q\to_*K(A,2)$ is merely null and (U) every $Q\to_*\Om K(B,4)$ is
merely null, then $q_*$ is bijective on classes. It is applied twice: to
$Q=P$ ((E) from (H4) and the defining null-homotopy of $P$; (U) is
$\hc P{K(B,3)}=0$) and to $Q=S^2$ ((E), (U): $\Om^2K(B,4)$ and $\Om^3K(B,4)$
are connected). The sphere case also produces $s$ (no choice: $s$ is the
unique class over a chosen class $\cls\sigma\in\hc{S^2}{K(\Z,2)}$), and a
commuting square of four maps of sets, three of them bijective, gives the
fourth.

Throughout, \emph{taking classes is never assumed to preserve fibres}:
existence of lifts is proved for maps, and injectivity of $q_*$ is proved
from a mere pointed equality of projections, through the uniqueness of
lifts up to mere equality (Lemma~\ref{l:uniq}). Lifts exist merely, and
any two lifts with equal pointed projections determine the same pointed
homotopy class (Theorem~\ref{t:lift}); connectedness of the full lift
spaces (pairs $(L,\alpha:L;q=F)$) is neither claimed nor needed.

\section{Conventions}

As in \code{doc/fibre-h3.pdf}: \code{isTrunc n X} is ``$X$ is
$(n-2)$-truncated'', \code{isConn n X} is ``$X$ is $(n-2)$-connected'';
$K_m=K(\Z,m)$, $K^A_m=K(A,m)$ (\code{EMMaps.KK}); $F;G$ is
\code{ptcomp F G}, with $(F;G)_1=G_1\circ F_1$ by definition; $0$ is
\code{ptconst}; $\cls F$ is \code{hcl F}, $a;b$ on classes is
\code{compC} with $\cls F;\cls G=\cls{F;G}$ (\code{comp-beta}); equality
of classes is mere equality of pointed maps (\code{effective},
\code{relates}). For a triple $(A,B,k)$, $k:K^A_2\to_*W$ with $W=K^B_4$,
$w_0=*$; $X=\Fib k=(y:K^A_2)\times(k\,y=w_0)$ pointed by $(*,k_0)$, and
$q=\mathrm{fpr}:X\to_*K^A_2$, $q(y,h)=y$, pointed by $\refl$.
$P=\Fib\kappa=(x:K_2)\times T_x$, $T_x=(\kappa x=*)$, $p_0=(*,\kappa_0)$,
$p=\mathrm{fpr}:P\to_*K_2$.

Levels used: $K^A_2$ is $\cC\,3$ and $\cT\,4$; $T_x$ is $\cC\,4$
(\code{conn-id}, inlined in \code{cp}); $W$ is $\cC\,5$, so $\Om W$ is $\cC\,4$, $\Om^2W$ is
$\cC\,3$, $\Om^3W$ is $\cC\,2$ (\code{conn-id}, \code{conn-down}); $X$ is
$\cC\,3$ (\code{Triple.conn3}), so $\Om X$ is $\cC\,2$.

\section{The lifting lemma}

Fix a pointed type $(Q,q_0)$ and a triple $(A,B,k)$. Write
$q_*(c)=c;\cls q$ for $c\in\hc QX$.

\begin{lemma}[pointing irrelevance; \code{FibreH3.pi-mere}]\label{l:pi}
If $\Om(Z,z_0)$ is connected, pointed maps $(k,\zeta),(k',\zeta')$ with
$k=k'$ are merely equal.
\end{lemma}

\begin{lemma}[lifts exist; \code{Fibres.Lift}]\label{l:ex}
If $\theta:F;k=0$ for $F:Q\to_*K^A_2$, there is $L:Q\to_*X$ with
$L;q=F$ (\code{lift}, \code{lift-fst}).
\end{lemma}

\begin{lemma}[lifts are merely unique, unpointed]\label{l:uniq0}
Assume (U$_0$): every $d:Q\to\Om W$ is merely equal to $\lambda\_.\refl$.
Let $u:Q\to X$, $g':Q\to K^A_2$, $h':(x:Q)\to k(g'x)=w_0$ and
$e:(\lambda x.\,(u\,x)_1)=g'$. Then merely $u=\lambda x.(g'x,h'x)$.
\end{lemma}
\begin{proof}
$J$ on $e$ ($g'$ and $h'$ free). For $g'\equiv\lambda x.(u\,x)_1$, let
$d\,x=(u\,x)_2^{-1}\cdot h'x:w_0=w_0$. By (U$_0$), merely
$\lambda\_.\refl=d$. Then for each $x$
\[(u\,x)_2\ \equiv\ (u\,x)_2\cdot\refl\ =\ (u\,x)_2\cdot d\,x\ =\ (u\,x)_2\cdot((u\,x)_2^{-1}\cdot h'x)\ =\ h'x\]
(\code{trans-cancel'}), so $u\,x=((u\,x)_1,h'x)$, and $u=\lambda
x.(g'x,h'x)$ by function extensionality.
\end{proof}

\begin{lemma}[(U$_0$) from (U)]\label{l:u0}
Assume (U): every $d:Q\to_*\Om W$ is merely equal to $0$. Then
(U$_0$) holds.
\end{lemma}
\begin{proof}
Let $d:Q\to\Om W$. Since $\Om W$ is connected, merely $r:d\,q_0=\refl$
(\code{conn-path}). By (U), merely $0=(d,r)$; its first component
gives $\lambda\_.\refl=d$.
\end{proof}

\begin{lemma}[lifts are merely unique]\label{l:uniq}
Assume (U). If $L,L':Q\to_*X$ and $L;q=L';q$, then merely $L=L'$.
\end{lemma}
\begin{proof}
The first component of $L;q=L';q$ is $e:(\lambda x.(L_1x)_1)=(\lambda
x.(L'_1x)_1)$ (by definition of \code{ptcomp} and \code{fpr}).
Lemma~\ref{l:uniq0} with $u=L_1$, $g'=\lambda x.(L'_1x)_1$,
$h'=\lambda x.(L'_1x)_2$ gives merely $L_1=L'_1$ ($\Sigma$-$\eta$);
Lemma~\ref{l:pi} ($\Om X$ connected) gives merely $L=L'$.
\end{proof}

\begin{theorem}[the lifting lemma]\label{t:lift}
Assume (E): for every $F:Q\to_*K^A_2$, merely $F;k=0$; and (U). Then
$q_*:\hc QX\to\hc Q{K^A_2}$ is a bijection.
\end{theorem}
\begin{proof}
\code{Bijection.Bij}: the target is a set; it suffices that $q_*$ is
injective and that each fibre is merely inhabited.
\emph{Surjective.} By induction on classes (into propositions), for
$\cls F$: by (E) merely $\theta:F;k=0$; Lemma~\ref{l:ex} gives $L$ with
$q_*\cls L=\cls{L;q}=\cls F$.
\emph{Injective.} By induction on both classes (the statement is a
proposition: classes form a set), for $\cls L,\cls{L'}$ with
$\cls{L;q}=\cls{L';q}$: merely $L;q=L';q$ (\code{effective}), so merely
$L=L'$ (Lemma~\ref{l:uniq}) and $\cls L=\cls{L'}$ (\code{relates}).
\end{proof}

So the lift space over each $F$ (pairs $(L,\,L;q=F)$) is merely
inhabited, and any two of its points have merely equal first components;
nothing is claimed about its higher structure.

\section{The source $P$}

\begin{lemma}[fibres of a first projection]\label{l:fibfst}
For $B:A\to U$ and $a:A$, the fibre of $\mathrm{pr}_1:\Sigma_AB\to A$
over $a$ is equivalent to $B\,a$, by $((a',b),\pi)\mapsto\pi_*b$ with
inverse $b\mapsto((a,b),\refl)$ ($J$ on $\pi$). Hence the fibres of
$p_1:P\to K_2$ are $\cC\,4$ (\code{equiv-conn}).
\end{lemma}

\begin{lemma}[$p^*$ on classes into $K^A_2$]\label{l:pstar}
$p^*:\hc{K_2}{K^A_2}\to\hc P{K^A_2}$, $c\mapsto\cls p;c$, is a bijection.
More precisely: for $F:P\to_*K^A_2$ let $G_F=(\bar F,\zeta_F)$ with
$\bar F=\mathrm{ext}\,F_1$ (\code{ConnectedMaps.Elim}, $m=4$, along
$p_1$, constant family $K^A_2$) and
$\zeta_F=(\bar F(*)=F_1p_0)\cdot F_0$ (\code{ext-beta} at $p_0$, since
$p_1p_0\equiv*$). Then merely $p;G_F=F$.
\end{lemma}
\begin{proof}
$(p;G_F)_1=\lambda w.\bar F(w_1)=F_1$ by \code{ext-beta} and function
extensionality; Lemma~\ref{l:pi} ($\Om K^A_2$ connected) gives merely
$p;G_F=F$. This is surjectivity (by induction on classes).
Injectivity: from $\cls{p;G}=\cls{p;G'}$, merely $p;G=p;G'$; the first
components give $G_1\circ p_1=G'_1\circ p_1$, so $G_1=G'_1$
(\code{Elim.unique}), and merely $G=G'$ by Lemma~\ref{l:pi}.
\code{Bij} concludes.
\end{proof}

\begin{lemma}[the fibre sequence]\label{l:fseq}
Merely $p;\kappa=0$.
\end{lemma}
\begin{proof}
$(p;\kappa)_1=\lambda w.\kappa(w_1)$ and $\lambda w.\,w_2$ is a homotopy
to $\lambda\_.*$; Lemma~\ref{l:pi} ($\Om K_4$ connected).
\end{proof}

\begin{proposition}[(E) for $P$]\label{p:EP}
For every triple $(A,B,k)$ and $F:P\to_*K^A_2$, merely $F;k=0$.
\end{proposition}
\begin{proof}
Let $G=G_F$ (Lemma~\ref{l:pstar}). By surjectivity of $\Phi_B$ (H4), for
$g=\Phi_B^{-1}\cls{G;k}$, $\cls{\kappa;K(g,4)}=\cls{G;k}$, so merely
$\kappa;K(g,4)=G;k$. Then, merely,
\[F;k=(p;G);k=p;(G;k)=p;(\kappa;K(g,4))=(p;\kappa);K(g,4)=0;K(g,4)=0\]
(Lemma~\ref{l:pstar}, \code{ptcomp-assoc}, Lemma~\ref{l:fseq},
\code{comp-const-l}).
\end{proof}

\begin{proposition}[(U) for $P$]\label{p:UP}
Every $d:P\to_*\Om W$ is merely $0$: \code{FibAt.merely-zero}
(\code{doc/fibre-h3.pdf}, Theorem~10, before the transfer along $e_3$).
\end{proposition}

\begin{corollary}\label{c:qP}
For every triple, $q_*:\hc P{\Fib k}\to\hc P{K^A_2}$ is a bijection.
\end{corollary}

\section{The source $S^2$}

Let $(S^2,\mathsf{base},\mathsf{surf})$ satisfy the sphere axiom and
$\mathrm{ev}=\code{evPt}:\Map(S^2,Z)\simeq\Om^2Z$ (\code{pointed-up}),
$\mathrm{ev}(F;G)=(\Om^2G)(\mathrm{ev}F)$ (\code{evPt-nat}).

\begin{lemma}\label{l:s2paths}
If $\Om^2(Z,z_0)$ is connected, any two $F,F':S^2\to_*Z$ are merely
equal: merely $\mathrm{ev}F=\mathrm{ev}F'$, and $\mathrm{ev}$ is
injective.
\end{lemma}

\begin{corollary}\label{c:qS}
(E) and (U) hold for $Q=S^2$ ($Z=W$: $\Om^2W$ is connected;
$Z=\Om W$: $\Om^3W$ is connected; for (E), $F;k$ and $0$ are two maps
$S^2\to_*W$). So $q_*:\hc{S^2}{\Fib k}\to\hc{S^2}{K^A_2}$ is a bijection
for every triple; in particular $p_*:\hc{S^2}P\to\hc{S^2}{K_2}$ is one
(the triple $(\Z,\Z,\kappa)$).
\end{corollary}

\begin{lemma}[looping is an equivalence; Corollary~11]\label{l:loopeq}
For $Y$ a 2-type, $\Om:\Map(K_2,Y)\to\Map(\Om K_2,\Om Y)$ is an
equivalence: every $\code{Unloop}\,f$ is contractible
(\code{Over.cor11}, $c=1$), which gives an inverse
$f\mapsto(\text{centre of Unloop}f)_1$, with $\Om$ of it $=f$ and, at
$f=\Om G$, $(G,\refl)$ equal to the centre.
\end{lemma}

\begin{lemma}[precomposition by a pointed equivalence]\label{l:preeq}
For $e:(X',x_0')\simeq_*(Z,z_0)$, $G\mapsto e_1;G$ is an equivalence
$\Map(Z,Y)\to\Map(X',Y)$ (pointed-equivalence induction,
\code{ptEqInd}; at the identity it is homotopic to the identity,
\code{ptcomp-idl}).
\end{lemma}

\begin{definition}[$\sigma$]\label{d:sigma}
Let $g:\Om K_1$ be the generator $\mathsf{unband}(1)$
(\code{CircleProperty.Circle.g}; it is locked, so the instances for
different parameters $Y$ are identified by a lemma, \code{g-def}: both
have $\mathsf{band}(g)=\mathsf{band}(g\cdot\refl)=1+\mathsf{band}(\refl)$
by \code{band-g}, then \code{band-inj}), $e_1:\Om K_2\simeq_*K_1$
(\code{CupSquare.e1}), $e_1^{-1}$ the map of its inverse \code{pt-inv e1}
(\code{CupSquare.e1inv}, by definition), and
\[c=(\Om e_1^{-1})(g)\ :\ \Om^2K_2,\qquad \sigma=\mathrm{ev}^{-1}(c):S^2\to_*K_2,\quad \mathrm{ev}\,\sigma=c .\]
\end{definition}

\begin{proposition}[$K_2$ is the 2-truncation of $S^2$, pointed]\label{p:sigma}
For every 2-type $(Y,y_0)$, $\sigma^*:\Map(K_2,Y)\to\Map(S^2,Y)$,
$G\mapsto\sigma;G$, is an equivalence.
\end{proposition}
\begin{proof}
Let $\Psi(G)=\mathrm{ev}^{\circ}(e_1^{-1};\Om G)$, where
$\mathrm{ev}^{\circ}:\Map(K_1,\Om Y)\simeq\Om^2Y$, $F\mapsto(\Om F)(g)$,
is the circle property (\code{Circle[$\Om Y$].pointed-up}). $\Psi$ is a
composite of three equivalences (Lemmas~\ref{l:loopeq}, \ref{l:preeq},
the circle property). And
\[\mathrm{ev}(\sigma;G)=(\Om^2G)(\mathrm{ev}\,\sigma)=(\Om^2G)((\Om e_1^{-1})(g))
=(\Om e_1^{-1};\Om^2G)(g)=(\Om(e_1^{-1};\Om G))(g)=\Psi(G)\]
(\code{evPt-nat}, Definition~\ref{d:sigma}, definition of \code{ptcomp},
\code{loop-comp}). So $\mathrm{ev}\circ\sigma^*=\Psi$ and $\sigma^*$ is an
equivalence (2-out-of-3, \code{cancel-equiv}, after \code{E.homotopy}:
the square commutes pointwise).
\end{proof}
Hence $\cls\sigma^*:\hc{K_2}Y\to\hc{S^2}Y$ is a bijection
(\code{pre-equiv}).

\begin{remark}[$p\circ s$ represents $1$]
Definition~\ref{d:sigma} makes $\sigma$ the generator: under
$\Om e_1:\Om^2K_2\simeq\Om K_1$ and the band $\Om K_1\simeq\Z$, $c$ goes
to $g$ and then to $1$. The proof uses only Proposition~\ref{p:sigma},
which is the property that characterises a generator; the value $1$ is not
used.
\end{remark}

\section{The mapping theorem for targets $\Fib k$}

\begin{definition}[$s$]\label{d:s}
$s=p_*^{-1}\cls\sigma\in\hc{S^2}P$ (Corollary~\ref{c:qS}), so
$s;\cls p=\cls\sigma$. A class, obtained without choice.
\end{definition}

\begin{proposition}\label{p:right}
For every abelian $A$, $s^*:\hc P{K^A_2}\to\hc{S^2}{K^A_2}$ is a bijection.
\end{proposition}
\begin{proof}
$s^*\circ p^*=\cls\sigma^*$: $s;(\cls p;d)=(s;\cls p);d=\cls\sigma;d$
(\code{comp-assoc}). $p^*$ is a bijection (Lemma~\ref{l:pstar}), and so
is $\cls\sigma^*$ (Proposition~\ref{p:sigma}, $K^A_2$ a 2-type); so
$s^*=\cls\sigma^*\circ(p^*)^{-1}$ is one (\code{cancel-r}: if $g\circ f$
and $f$ are equivalences, so is $g$).
\end{proof}

\begin{theorem}[the mapping theorem for $\Fib k$]\label{t:main}
For every triple $(A,B,k)$, $s^*:\hc P{\Fib k}\to\hc{S^2}{\Fib k}$ is a
bijection.
\end{theorem}
\begin{proof}
The square
\[\begin{array}{ccc}
\hc P{\Fib k} & \xrightarrow{\ q_*\ } & \hc P{K^A_2}\\
{\scriptstyle s^*}\downarrow & & \downarrow{\scriptstyle s^*}\\
\hc{S^2}{\Fib k} & \xrightarrow{\ q_*\ } & \hc{S^2}{K^A_2}
\end{array}\]
commutes by \code{comp-assoc}: $(s;c);\cls q=s;(c;\cls q)$. The top
(Corollary~\ref{c:qP}), the right (Proposition~\ref{p:right}) and the
bottom (Corollary~\ref{c:qS}) are bijections; so the left is one
(2-out-of-3, \code{cancel-equiv}).
\end{proof}

\begin{corollary}[transfer]\label{c:transfer}
If $\Fib k\simeq_*(X,x_0)$ for some triple, $s^*:\hc PX\to\hc{S^2}X$ is
a bijection (\code{ptEqInd}). Since the statement is a proposition, a
\emph{mere} pointed equivalence suffices.
\end{corollary}

\begin{remark}[what remains for \code{MapThm}]
\code{Pi3Sphere.Conditional.MapThm} asks for $s^*$ bijective at every
simply connected 3-type. \code{SphereUniversal.Assemble} uses it at
$X=P$ (Theorem~\ref{t:main} with the triple $(\Z,\Z,\kappa)$) and at
$X=T=\tr{S^2}3$; by Corollary~\ref{c:transfer}, the latter follows from a
mere Postnikov presentation $T\simeq_*\Fib k$ (\S3d, not done here).
\end{remark}

\section{Correspondence with the formal proof}\label{sec:corr}
\emph{Status: checked} (\code{./segc sag/pi3-s2-postnikov all}: 70 modules,
0 failed; a fresh process: 3003 declarations, 0 failed, 1.2\,s; every
\code{MapFib} declaration under 0.03\,s; \code{mapfib.opens}, 110 entries,
every line ``nothing else opened''). Module \code{MapFib}:
\code{fib-fst-ret}, \code{fib-fst} (Lemma~\ref{l:fibfst});
\code{cancel-snd} (the step of Lemma~\ref{l:uniq0}); \code{cancel-r};
\code{LevW[B]}: \code{cOW}, \code{cOOW}, \code{cOOOW} (levels);
\code{Lifting[Q,q0,A,B,k]}: \code{Ex}, \code{Un}, \code{U0} ((E), (U),
(U$_0$)), \code{qs} ($q_*$), \code{cOF}, \code{u0} (Lemma~\ref{l:u0}),
\code{uniq0} (Lemma~\ref{l:uniq0}), \code{uniq} (Lemma~\ref{l:uniq}),
\code{Main[ex,un]}: \code{surj}, \code{inj-at}, \code{inj}, \code{equiv}
(Theorem~\ref{t:lift}); \code{SourceP[$\kappa$,h4k]}: \code{cp}
(Lemma~\ref{l:fibfst}), \code{PStar[A]}: \code{EA}, \code{cOA},
\code{ext-pt} ($G_F$), \code{ext-eq}, \code{pstar}, \code{surj},
\code{inj-at}, \code{inj}, \code{equiv} (Lemma~\ref{l:pstar}); \code{cOK4},
\code{fseq} (Lemma~\ref{l:fseq}); \code{ExP[A,B,k,F]}: \code{g},
\code{g-eq}, \code{chain}, \code{null}; \code{ex-P}
(Proposition~\ref{p:EP}); \code{un-P} (Proposition~\ref{p:UP}); \code{qP}
(Corollary~\ref{c:qP}); \code{pre-pteq} (Lemma~\ref{l:preeq});
\code{SourceS[S2,base,surf,ump]}: \code{s2-paths} (Lemma~\ref{l:s2paths}),
\code{ex-S}, \code{un-S}, \code{qS} (Corollary~\ref{c:qS}),
\code{LoopEq[Y,y0,tY]}: \code{cell}, \code{equiv} (Lemma~\ref{l:loopeq}),
\code{gen-c}, \code{sigma}, \code{sigma-ev}, \code{g-def}
(Definition~\ref{d:sigma}), \code{Sig[Y,y0,tY]}: \code{Psi},
\code{psi-equiv}, \code{sigma-nat}, \code{sigma-star}, \code{sigma-cls}
(Proposition~\ref{p:sigma}); \code{Final[S2,base,surf,ump,$\kappa$,h4k]}:
\code{s}, \code{s-eq} (Definition~\ref{d:s}), \code{right}
(Proposition~\ref{p:right}), \code{rep-fib} (Theorem~\ref{t:main}),
\code{rep-pteq} (Corollary~\ref{c:transfer} for a pointed equivalence; the
mere version, \code{TC.rec} into the proposition \code{isEquiv}, is not
written); \code{rep-cup}: Theorem~\ref{t:main} for \code{CupSquare.kappa},
\code{H4Proof.h4}.
\end{document}
